\chapter{Special and General Relativity}\label{general-relativity}
\chapterstatus{planned}

Relativity replaces Galilean spacetime by a Lorentzian manifold whose curvature is gravity. The
mathematics is pseudo-Riemannian geometry and hyperbolic partial differential equations; this chapter
develops it from special relativity to black holes and the singularity theorems.

\noindent\textbf{Prerequisites.} Chapter~\ref{smooth-manifolds} (Smooth Manifolds), Chapter~\ref{connections} (Connections and Curvature), Chapter~\ref{riemannian-geometry} (Riemannian Geometry), Chapter~\ref{sobolev-spaces} (Partial Differential Equations) and Chapter~\ref{classical-mechanics} (Classical Mechanics).

\noindent\textbf{Planned contents.}
\begin{enumerate}
\item Special relativity: Minkowski space, the Lorentz and Poincar\'e groups, causal structure,
proper time, the twin paradox, relativistic energy and momentum; Maxwell's equations in tensor form.
\item Lorentzian geometry: causal vectors and curves, time orientation, geodesics and the
Levi-Civita connection in Lorentzian signature, the reversed triangle inequality.
\item Einstein's equations: the stress--energy tensor, the Einstein--Hilbert action, the vacuum
equations; the Newtonian limit.
\item The Schwarzschild solution: Birkhoff's theorem, the perihelion of Mercury, the bending of light,
Kruskal coordinates and the black hole; Kerr spacetime (survey).
\item Cosmology: Friedmann--Lema\^itre--Robertson--Walker models, the Big Bang.
\item Causality theory: global hyperbolicity, Cauchy surfaces, Geroch's theorem; the Cauchy problem
for Einstein's equations (Choquet-Bruhat, statement).
\item The singularity theorems of Penrose and Hawking; the positive mass theorem (statement).
\item Gravitational waves: the linearised equations.
\item Exercises with solutions.
\end{enumerate}
